% concatenated from main.tex
 \documentclass[11pt]{amsart}

\usepackage[a4paper,margin=3.2cm]{geometry}
\usepackage{amsmath,amssymb,amsthm,mathtools}
\usepackage[T1]{fontenc}
\usepackage[utf8]{inputenc}
\usepackage{lmodern}
\usepackage{hyperref}

\hypersetup{
  colorlinks=true,
  linkcolor=blue,
  citecolor=blue,
  urlcolor=blue
}

\newtheorem{theorem}{Theorem}
\newtheorem{proposition}[theorem]{Proposition}
\newtheorem{lemma}[theorem]{Lemma}
\newtheorem{corollary}[theorem]{Corollary}
\theoremstyle{definition}
\newtheorem{definition}[theorem]{Definition}
\theoremstyle{remark}
\newtheorem{remark}[theorem]{Remark}

\newcommand{\C}{\mathbb C}
\newcommand{\R}{\mathbb R}
\newcommand{\D}{\mathbb D}
\newcommand{\re}{\operatorname{Re}}
\newcommand{\im}{\operatorname{Im}}
\newcommand{\rank}{\operatorname{rank}}
\newcommand{\id}{\operatorname{id}}

\title[A Lewy theorem for pluriharmonic maps in $\C^2$]
{A Lewy-type theorem for pluriharmonic mappings in $\C^2$ and bi-Lipschitz quasiconformal harmonic maps}

\author{David Kalaj}
\address{University of Montenegro, Faculty of Natural Sciences and Mathematics, 81000 Podgorica, Montenegro}
\email{davidk@ucg.ac.me}

\subjclass[2020]{Primary 32U05; Secondary 30C65, 32S55, 31B05, 57K10}
\keywords{Pluriharmonic mapping, Lewy theorem, Jacobian, Milnor fibration, plane curve singularity, local flatness, quasiconformal mapping, Dini-smooth domain}

\begin{document}


\begin{abstract}
Lewy's theorem says that a one-to-one harmonic mapping between plane
domains has nonvanishing Jacobian. In dimensions at least three this
statement is false for general harmonic homeomorphisms. We prove a
four-dimensional Lewy theorem under the additional complex-analytic
assumption of pluriharmonicity. More precisely, if
\[
F:\Omega\subset \mathbb C^2\to \mathbb C^2
\]
is a \(C^2\) pluriharmonic mapping which is one-to-one in a neighborhood
of a point \(p\), then the real Jacobian of \(F\) is nonzero at \(p\).

The proof is local. If the Jacobian vanished, a nontrivial real linear
projection of \(F\) would be the real part of a holomorphic function
\(f\) with \(df(p)=0\). The level hypersurface of this projection would
therefore be a real analytic germ of the form
\[
\{\operatorname{Re} f=0\}.
\]
We prove that such a germ cannot be locally flat at a critical point of
\(f\). The obstruction is topological: local flatness forces the local
homology of the hypersurface germ to agree with that of a real
hyperplane, and hence forces every sufficiently small admissible link
to have the integral homology, in particular the Euler characteristic,
of \(S^2\). In the reduced case the Milnor open book of the plane curve
singularity \(f^{-1}(0)\) gives a contradictory Euler-characteristic
formula, while in the nonreduced case the local normal form produces
more than two local complementary components.

We then prove a separate bi-Lipschitz criterion for harmonic
quasiconformal mappings: a harmonic quasiconformal homeomorphism from
the unit ball onto a bounded \(C^1\)-Dini domain is bi-Lipschitz,
provided it is already a local \(C^1\)-diffeomorphism. This part uses a
regularized defining function, the quasihyperbolic metric, and a
blow-up argument whose limiting half-space height function is forced to
be linear. The Lewy-type theorem is used only afterward, to verify the
local-diffeomorphism hypothesis for quasiconformal pluriharmonic
homeomorphisms in \(\mathbb C^2\). In the planar case, Lewy's classical
theorem supplies the corresponding local nondegeneracy for harmonic
quasiconformal homeomorphisms between bounded finitely connected
\(C^1\)-Dini domains.
\end{abstract}


\maketitle

\section{Introduction}

Lewy's classical theorem \cite{Lewy1936} asserts that if
\[
        F=(u,v):U\subset\R^2\longrightarrow\R^2
\]
is harmonic and one-to-one, then its Jacobian determinant does not
vanish. Thus a univalent planar harmonic mapping is locally a
diffeomorphism. Lewy later proved a related result for harmonic
gradients in dimension three \cite{Lewy1968}. The unrestricted
higher-dimensional analogue, however, is false: Wood \cite{Wood1991}
constructed a harmonic homeomorphism of \(\R^3\) whose Jacobian
vanishes on a plane, and the same construction gives counterexamples in
every real dimension \(n\ge3\).

The purpose of this paper is to show that a positive four-dimensional
result is recovered if ordinary harmonicity is replaced by
pluriharmonicity. A complex-valued \(C^2\) function is pluriharmonic if
\[
        \partial\bar\partial f=0,
\]
and a mapping \(F:\Omega\subset\C^2\to\C^2\) is pluriharmonic if each
of its complex components is pluriharmonic. Equivalently, after the
identification \(\C^2\cong\R^4\), all four real coordinate functions of
\(F\) are pluriharmonic.

Naser \cite{Naser1973} claimed a Lewy-type assertion for
pluriharmonic mappings in real dimension four. The argument in
\cite{Naser1973} contains an important local-topological point which is
not justified as stated. More precisely, the proof first obtains a
component count for the complement of the set $\{h_1=0\}$ by looking at
one-dimensional complex slices. It then passes to the inverse image of a
sufficiently small sphere under the homeomorphism $h$ and compares this
with the fact that a real hyperplane cuts a three-sphere into two
components. This passage does not by itself prove the required local
statement: the inverse image of a small sphere is a link of the germ,
not a neighborhood of the origin, and a component count in complex-line
slices need not determine the topology of the complement in
$\mathbb R^4$.

A simple model illustrates the difficulty. Let
$f(z_1,z_2)=z_1^2+z_2^2$. On the complex line $\{z_2=0\}$, the equation
$\operatorname{Re} f=0$ becomes $\operatorname{Re} z_1^2=0$, and hence
the punctured disc is divided into four sectors. However, in
$\mathbb R^4\cong\mathbb C^2$ one has
\[
        \operatorname{Re}(z_1^2+z_2^2)
        =
        x_1^2+x_2^2-y_1^2-y_2^2,
        \qquad z_j=x_j+iy_j.
\]
Thus the complement of $\{\operatorname{Re} f=0\}$ near the origin has
the two regions
\[
        x_1^2+x_2^2>y_1^2+y_2^2,
        \qquad
        x_1^2+x_2^2<y_1^2+y_2^2.
\]
Hence the number of sectors seen in a complex-line slice does not
determine the number of local complementary components in
$\mathbb R^4$. Moreover, on a small sphere $S^3_\varepsilon$ the link is
\[
        S^3_\varepsilon\cap\{\operatorname{Re}(z_1^2+z_2^2)=0\}
        =
        \{|x|=|y|=\varepsilon/\sqrt 2\},
\]
which is a two-torus, not a two-sphere. This example shows that the
correct invariant is the topology of the embedded link, rather than a
component count in one-dimensional slices.

The present paper supplies this missing local argument in complex
dimension two. We show that if $f$ is holomorphic near $0\in\mathbb C^2$
and $df(0)=0$, then the real hypersurface germ
$\{\operatorname{Re} f=0\}$ is not locally flat at $0$. We use the term \emph{locally flat} in the topological sense for
embedded germs. Thus an embedded hypersurface germ \((\R^4,X,0)\) is
locally flat at \(0\) if there are neighborhoods \(U,V\) of \(0\) in
\(\R^4\), a real hyperplane \(H\subset\R^4\), and a homeomorphism
\[
        \Phi:U\to V
\]
such that
\[
        \Phi(0)=0,
        \qquad
        \Phi(X\cap U)=H\cap V.
\]
In other words, after a local topological change of coordinates in the
ambient space, the set \(X\) becomes a real hyperplane. The obstruction
is detected by the local homology of the hypersurface germ and, in the
reduced case, by the Euler characteristic of the link
$S^3_\varepsilon\cap\{\operatorname{Re} f=0\}$ as computed from
Milnor's fibration theorem for the plane curve singularity $f^{-1}(0)$.
This replaces the component-counting argument in complex-line slices
by an invariant of the embedded germ that is directly preserved by
local flatness.


\begin{theorem}\label{th:main}
Let \(\Omega\subset\C^2\) be a domain, and let
\[
        F:\Omega\longrightarrow\C^2
\]
be a \(C^2\) pluriharmonic mapping. If \(F\) is one-to-one in a
neighborhood of a point \(p\in\Omega\), then the real Jacobian of \(F\)
does not vanish at \(p\):
\[
        J_F(p)\ne0.
\]
\end{theorem}

The theorem is local both in the source and in the target. Thus global
univalence may be replaced by injectivity in a sufficiently small
neighborhood of the point under consideration. Throughout the paper,
``Jacobian'' means the real Jacobian determinant of the corresponding
mapping \(F:\R^4\to\R^4\).

The proof has two independent ingredients. The analytic ingredient is
simple: every real-valued pluriharmonic function is locally the real
part of a holomorphic function. Thus, if \(J_F(p)=0\), a nonzero real
linear projection of \(F\) gives a holomorphic function \(f\) satisfying
\(df(p)=0\). The geometric ingredient is the following topological
obstruction: for a nonconstant holomorphic function \(f\) in two
variables, the germ
\[
        \{\operatorname{Re} f=0\}
\]
cannot be locally flat at a critical point of \(f\). In the reduced
case this is detected by the local homology of the germ together with
the Euler characteristic of the Milnor open-book pages; in the
nonreduced case it follows from the local product normal form and the
number of local complementary components.

After proving Theorem~\ref{th:main}, we turn to quasiconformal
harmonic mappings. The bi-Lipschitz argument separates two issues.
First, boundary regularity of the image domain gives two-sided estimates
for the distance of \(F(z)\) to the boundary and yields a Lipschitz
bound. Second, a Lewy-type nondegeneracy theorem prevents the minimal
stretch from degenerating in the interior. This separation gives a
criterion which applies to any harmonic quasiconformal class satisfying
a Lewy property and stable under blow-ups.

Applying this criterion to the pluriharmonic class in \(\C^2\), using
Theorem~\ref{th:main}, gives the following consequence.

\begin{theorem}\label{th:intro-bilip}
Let \(\Omega\subset\C^2\cong\R^4\) be a bounded domain with
\(C^1\)-Dini boundary, and let
\[
        F:\mathbb B^2\to\Omega
\]
be a \(K\)-quasiconformal pluriharmonic homeomorphism. Then \(F\) is
bi-Lipschitz.
\end{theorem}

In the plane, the same method uses Lewy's classical theorem instead of
Theorem~\ref{th:main}. It gives the following more flexible form, in
which both the source and the image may be finitely connected Dini-smooth
domains.

\begin{theorem}\label{th:intro-planar-bilip}
Let \(D,\Omega\subset\C\) be bounded finitely connected domains whose
boundary components are \(C^1\)-Dini Jordan curves. Let
\[
        f:D\to\Omega
\]
be a harmonic \(K\)-quasiconformal homeomorphism. Then \(f\) is
bi-Lipschitz.
\end{theorem}

The proof of these bi-Lipschitz results is given in the final section.
It uses a regularized distance function for \(C^1\)-Dini domains,
quasihyperbolic distortion estimates, compactness of quasiconformal
mappings, and the appropriate Lewy theorem.

The paper is organized as follows. Section~2 collects the analytic
preliminaries used in the local reduction: the primitive lemma for
real-valued pluriharmonic functions and the standard local factorization
facts for plane curve germs. Section~3 proves the topological obstruction to
local flatness of \(\{\operatorname{Re} f=0\}\). Section~4 proves the
Lewy theorem. Section~5 proves the bi-Lipschitz theorem for harmonic
quasiconformal local diffeomorphisms and then applies it to pluriharmonic
mappings in \(\C^2\) and to planar harmonic mappings between Dini-smooth
domains.

\section{Analytic preliminaries}

We collect two standard local facts used in the proof of the Lewy-type
theorem: the existence of holomorphic primitives for real-valued
pluriharmonic functions, and the local factorization of plane curve germs.
We begin with the elementary primitive result used in the proof of the main
theorem.

\begin{lemma}\label{lem:primitive}
Let $U\subset\C^2$ be a simply connected domain, and let
$u\in C^2(U,\R)$ be pluriharmonic. Then there is a holomorphic function
$f$ on $U$ such that
\[
        u=\re f.
\]
The function $f$ is unique up to addition of a purely imaginary
constant.
\end{lemma}

\begin{proof}
Since $u$ is pluriharmonic,
\[
        \bar\partial\partial u=0.
\]
Set
\[
        \omega=2\partial u.
\]
Then $\omega$ is a holomorphic $(1,0)$-form. Moreover,
\[
        d\omega=\partial\omega+\bar\partial\omega=0,
\]
because $\partial^2=0$ and $\bar\partial\partial u=0$. Since $U$ is
simply connected, the closed holomorphic one-form $\omega$ has a
holomorphic primitive. Hence there exists a holomorphic function $f$
such that
\[
        df=\omega=2\partial u.
\]
Taking real parts gives
\[
        d(\re f)=\re(df)=\re(2\partial u)=du.
\]
Thus $u-\re f$ is constant. Adding a real constant to $f$, we may
arrange that $u=\re f$. If $f$ and $g$ are two such holomorphic
functions, then $\re(f-g)=0$, so $f-g$ is constant and purely imaginary.
\end{proof}

\subsection{Local factorization of plane curve germs}

We shall use standard facts from local analytic geometry and plane curve
singularities. References include Grauert--Remmert
\cite[Ch.~2]{GrauertRemmert1984}, H\"ormander \cite[Ch.~VI]{Hormander1990},
Chirka \cite{Chirka1989}, Wall \cite[Ch.~2]{Wall2004}, and
Brieskorn--Kn\"orrer \cite[Ch.~8]{BrieskornKnorrer1986}.

\begin{definition}\label{def:convergent-power-series}
The ring
\[
        \C\{z_1,z_2\}
\]
is the local ring of convergent power series at $0\in\C^2$. Thus an
element of $\C\{z_1,z_2\}$ is a power series
\[
        \sum_{j,k\ge0} a_{jk}z_1^jz_2^k
\]
which converges in some neighborhood of the origin, considered as a germ
at $0$.
\end{definition}

\begin{definition}\label{def:unit-irreducible-associated}
Let $R=\C\{z_1,z_2\}$.
\begin{enumerate}
\item A function $u\in R$ is a \emph{unit} if it has a multiplicative
inverse in $R$. Equivalently, $u(0)\ne0$. After shrinking the
representative neighborhood, this means that $u$ is nowhere zero.

\item A nonzero nonunit $g\in R$ is \emph{irreducible} if every
factorization
\[
        g=ab,
        \qquad a,b\in R,
\]
forces either $a$ or $b$ to be a unit.

\item Two nonzero elements $g,h\in R$ are \emph{associated} if
\[
        g=vh
\]
for some unit $v\in R$. Otherwise they are \emph{nonassociated}.
\end{enumerate}
\end{definition}

Geometrically, an irreducible factor in $\C\{z_1,z_2\}$ defines one
local branch of a plane analytic curve. Associated irreducible factors
define the same branch, while pairwise nonassociated irreducible factors
define distinct branches.

\begin{proposition}\label{prop:ufd-convergent-power-series}
The local analytic ring
\[
        \C\{z_1,z_2\}
\]
is a unique factorization domain.
\end{proposition}

\begin{proof}
The ring $\C\{z_1,z_2\}$ is a regular local ring. Every regular local
ring is factorial. Hence $\C\{z_1,z_2\}$ is a unique factorization
domain. See, for example, Matsumura \cite[Theorem~20.3]{Matsumura1986}.
\end{proof}

\begin{theorem}[Local factorization theorem]\label{th:local-factorization}
Let $f$ be a nonzero holomorphic function near $0\in\C^2$, and suppose
that
\[
        f(0)=0.
\]
After shrinking the neighborhood of $0$, one can write
\[
        f=u f_1^{m_1}\cdots f_r^{m_r},
\]
where $u$ is holomorphic and nowhere zero, the functions
$f_1,\ldots,f_r$ are irreducible elements of $\C\{z_1,z_2\}$, no two of
them are associated, and $m_j\ge1$. The germs $\{f_j=0\}$ are precisely
the distinct irreducible local branches of the plane curve germ
$\{f=0\}$, and $m_j$ is the multiplicity of the corresponding branch.
\end{theorem}

\begin{proof}
The Taylor series of $f$ defines an element of $\C\{z_1,z_2\}$. By
Proposition~\ref{prop:ufd-convergent-power-series}, this ring is a
unique factorization domain. Hence $f$ has a factorization
\[
        f=u f_1^{m_1}\cdots f_r^{m_r},
\]
where $u$ is a unit, the $f_j$ are pairwise nonassociated irreducible
elements, and $m_j\ge1$. Since $u(0)\ne0$, after shrinking the
representative neighborhood the function $u$ is nowhere zero. Thus $u$
does not affect the zero set. The irreducible factors $f_j$ give
exactly the distinct local analytic branches of $\{f=0\}$, and the
exponent $m_j$ records the multiplicity of the corresponding branch.
\end{proof}

\begin{definition}\label{def:reduced-plane-curve}
The germ $f\in\C\{z_1,z_2\}$ is called \emph{reduced} if in the
factorization
\[
        f=u f_1^{m_1}\cdots f_r^{m_r}
\]
all multiplicities satisfy $m_j=1$.
\end{definition}

\begin{remark}\label{rem:factorization-examples}
For example, $z_1$ is irreducible and defines the smooth branch
$\{z_1=0\}$. The function
\[
        z_1^2-z_2^2=(z_1-z_2)(z_1+z_2)
\]
has two distinct branches. By contrast, $z_1-z_2$ and $2(z_1-z_2)$ are
associated and define the same branch.
\end{remark}


\section{A topological obstruction to local flatness}

The central point of the proof is the following proposition. It says
that the real hypersurface \(\{\re f=0\}\) cannot be locally flat at a
critical point of the holomorphic function \(f\).

We use local flatness in the sense of embedded germs. Thus the germ
\((\R^4,X,0)\) is locally homeomorphic to the germ of a real hyperplane if
there are neighborhoods \(U,V\) of the origins in \(\R^4\), a real
hyperplane \(H\subset\R^4\), and a homeomorphism of pairs
\[
        (U,X\cap U,0)\longrightarrow (V,H\cap V,0).
\]

We first record the elementary consequence of local flatness which will
be used below. The point of the lemma is that local flatness determines
the local homology of the hypersurface germ and therefore the reduced
homology of every sufficiently small admissible link.


Throughout this section, \(H_k(\,\cdot\,;\mathbb Z)\) denotes singular
homology with integer coefficients, \(\widetilde H_k(\,\cdot\,;\mathbb Z)\)
denotes reduced singular homology, and \(H_k(A,B;\mathbb Z)\) denotes
relative singular homology of the pair \((A,B)\). We write
\(B_\varepsilon\subset\mathbb R^4\) for the closed Euclidean ball of radius
\(\varepsilon\) centered at the origin and
\(S^3_\varepsilon=\partial B_\varepsilon\) for its boundary.

Recall that, for a nonempty connected space \(Y\),
\[
\widetilde H_0(Y;\mathbb Z)=0,
\]
while
\[
\widetilde H_k(Y;\mathbb Z)=H_k(Y;\mathbb Z),
\qquad k>0.
\]
We shall also use the standard fact that, for a topological space \(L\), its
cone
\[
\operatorname{Cone}(L)
=
(L\times[0,1])/(L\times\{0\})
\]
is contractible, and that after removing the cone vertex \(v\), the punctured
cone
\[
\operatorname{Cone}(L)\setminus\{v\}
\]
deformation retracts onto \(L\).

\begin{lemma}\label{lem:flat-link-homology}
Let \(X\subset\mathbb R^4\) be a real analytic hypersurface germ at \(0\).
Assume that the embedded germ \((\mathbb R^4,X,0)\) is locally flat; that is,
it is homeomorphic, as a germ of pairs, to the germ
\[
(\mathbb R^4,\mathbb R^3\times\{0\},0).
\]
Then, for every sufficiently small admissible \(\varepsilon>0\), the link
\[
L_\varepsilon=X\cap S^3_\varepsilon
\]
has the integral homology of \(S^2\). More precisely,
\[
\widetilde H_j(L_\varepsilon;\mathbb Z)
\cong
\begin{cases}
\mathbb Z, & j=2,\\
0, & j\ne2.
\end{cases}
\]
In particular,
\[
\chi(L_\varepsilon)=2.
\]
\end{lemma}

\begin{proof}
Choose \(\varepsilon>0\) sufficiently small so that both the local-flatness
description and the local conical structure of \(X\) are valid in
\(B_\varepsilon\).

Since the embedded germ \((\mathbb R^4,X,0)\) is locally flat, a sufficiently
small neighborhood of \(0\) in \(X\) is homeomorphic to a sufficiently small
neighborhood of \(0\) in the hyperplane
\[
\mathbb R^3\times\{0\}\subset\mathbb R^4.
\]
Hence the local homology groups of \(X\) at \(0\) coincide with those of
\(\mathbb R^3\) at the origin. More precisely,
\[
H_k\bigl(
X\cap B_\varepsilon,
(X\cap B_\varepsilon)\setminus\{0\};
\mathbb Z
\bigr)
\cong
H_k\bigl(
\mathbb R^3,
\mathbb R^3\setminus\{0\};
\mathbb Z
\bigr).
\]

We recall briefly why the relative homology groups on the right can be
computed from the homology of the punctured space. For every pair of
topological spaces \(B\subset A\), the relative singular chain group is
defined by
\[
C_k(A,B;\mathbb Z)
=
C_k(A;\mathbb Z)/C_k(B;\mathbb Z).
\]
Thus there is a short exact sequence of chain complexes
\[
0
\longrightarrow
C_\bullet(B;\mathbb Z)
\longrightarrow
C_\bullet(A;\mathbb Z)
\longrightarrow
C_\bullet(A,B;\mathbb Z)
\longrightarrow
0.
\]
The corresponding long exact sequence in homology is
\[
\cdots
\longrightarrow
H_k(B;\mathbb Z)
\longrightarrow
H_k(A;\mathbb Z)
\longrightarrow
H_k(A,B;\mathbb Z)
\overset{\partial}{\longrightarrow}
H_{k-1}(B;\mathbb Z)
\longrightarrow
H_{k-1}(A;\mathbb Z)
\longrightarrow
\cdots .
\]
Here \(\partial\) denotes the connecting homomorphism. The exactness of this
sequence means that, at each term, the image of the preceding map coincides
with the kernel of the following map.

In particular, if \(A\) is contractible, then its reduced homology groups
vanish. The long exact sequence therefore gives the standard degree-shift
isomorphism
\[
H_k(A,B;\mathbb Z)
\cong
\widetilde H_{k-1}(B;\mathbb Z).
\]

We now apply this observation to
\[
A=\mathbb R^3,
\qquad
B=\mathbb R^3\setminus\{0\}.
\]
The space \(\mathbb R^3\) is contractible, whereas
\(\mathbb R^3\setminus\{0\}\) deformation retracts radially onto \(S^2\).
Indeed, the radial projection
\[
r:\mathbb R^3\setminus\{0\}\longrightarrow S^2,
\qquad
r(x)=\frac{x}{\|x\|},
\]
is part of a deformation retraction. Consequently,
\[
\mathbb R^3\setminus\{0\}\simeq S^2,
\]
where \(\simeq\) denotes homotopy equivalence.

Hence
\[
H_k\bigl(
\mathbb R^3,
\mathbb R^3\setminus\{0\};
\mathbb Z
\bigr)
\cong
\widetilde H_{k-1}\bigl(
\mathbb R^3\setminus\{0\};
\mathbb Z
\bigr)
\cong
\widetilde H_{k-1}(S^2;\mathbb Z).
\]
Since the reduced integral homology of the \(2\)-sphere is
\[
\widetilde H_j(S^2;\mathbb Z)
\cong
\begin{cases}
\mathbb Z, & j=2,\\
0, & j\ne2,
\end{cases}
\]
the only nonzero relative homology group occurs when
\[
k-1=2,
\]
that is, when
\[
k=3.
\]
Therefore
\[
H_k\bigl(
X\cap B_\varepsilon,
(X\cap B_\varepsilon)\setminus\{0\};
\mathbb Z
\bigr)
\cong
\begin{cases}
\mathbb Z, & k=3,\\
0, & k\ne3.
\end{cases}
\tag{3.1}
\]

We next relate these local homology groups to the topology of the link
\(L_\varepsilon\). Choose a Whitney stratification of \(X\) for which
\(\{0\}\) is a stratum. By the local conical structure theorem for real
analytic, and more generally subanalytic, sets, for every sufficiently small
admissible \(\varepsilon>0\) there is a homeomorphism of pairs
\[
\bigl(
X\cap B_\varepsilon,
(X\cap B_\varepsilon)\setminus\{0\}
\bigr)
\cong
\bigl(
\operatorname{Cone}(L_\varepsilon),
\operatorname{Cone}(L_\varepsilon)\setminus\{v\}
\bigr),
\]
where
\[
L_\varepsilon=X\cap S^3_\varepsilon
\]
and \(v\) denotes the vertex of the cone; see
\cite{Lojasiewicz1964},
\cite[Sec.~9.3]{BochnakCosteRoy1998}, and
\cite{BurgheleaVerona1972}.

Thus, topologically, a sufficiently small neighborhood of the origin in
\(X\) is obtained by taking the link \(L_\varepsilon\) and forming its cone,
with the origin corresponding to the cone vertex.

Set
\[
C=\operatorname{Cone}(L_\varepsilon)
\]
and
\[
C^\ast=C\setminus\{v\}.
\]
The cone \(C\) is contractible, and therefore
\[
\widetilde H_j(C;\mathbb Z)=0
\]
for every \(j\). Moreover,
\[
C^\ast\cong L_\varepsilon\times(0,1],
\]
and hence \(C^\ast\) deformation retracts onto
\[
L_\varepsilon\times\{1\}\cong L_\varepsilon.
\]
Consequently,
\[
C^\ast\simeq L_\varepsilon.
\]

Applying the preceding general observation to the pair \((C,C^\ast)\), and
using the contractibility of \(C\), we obtain
\[
H_k(C,C^\ast;\mathbb Z)
\cong
\widetilde H_{k-1}(C^\ast;\mathbb Z).
\]
Since \(C^\ast\simeq L_\varepsilon\), homotopy invariance of singular
homology gives
\[
\widetilde H_{k-1}(C^\ast;\mathbb Z)
\cong
\widetilde H_{k-1}(L_\varepsilon;\mathbb Z).
\]
Therefore
\[
H_k\bigl(
\operatorname{Cone}(L_\varepsilon),
\operatorname{Cone}(L_\varepsilon)\setminus\{v\};
\mathbb Z
\bigr)
\cong
\widetilde H_{k-1}(L_\varepsilon;\mathbb Z).
\tag{3.2}
\]

Combining the local-flatness computation in \((3.1)\) with the conical
identification in \((3.2)\), we obtain
\[
\widetilde H_{k-1}(L_\varepsilon;\mathbb Z)
\cong
\begin{cases}
\mathbb Z, & k=3,\\
0, & k\ne3.
\end{cases}
\]
Putting
\[
j=k-1,
\]
we conclude that
\[
\widetilde H_j(L_\varepsilon;\mathbb Z)
\cong
\begin{cases}
\mathbb Z, & j=2,\\
0, & j\ne2.
\end{cases}
\]

Thus \(L_\varepsilon\) has the integral homology of \(S^2\). Equivalently,
its ordinary homology groups satisfy
\[
H_0(L_\varepsilon;\mathbb Z)\cong\mathbb Z,
\]
\[
H_1(L_\varepsilon;\mathbb Z)=0,
\]
and
\[
H_2(L_\varepsilon;\mathbb Z)\cong\mathbb Z,
\]
while
\[
H_j(L_\varepsilon;\mathbb Z)=0,
\qquad j\ge3.
\]

Finally, \(L_\varepsilon\) is a compact subanalytic set. Hence it is
triangulable and has the homotopy type of a finite CW complex. Its Euler
characteristic can therefore be computed from its homology by
\[
\chi(L_\varepsilon)
=
\sum_{j\ge0}
(-1)^j
\operatorname{rank}
H_j(L_\varepsilon;\mathbb Z).
\]
Using the homology groups above, we obtain
\[
\chi(L_\varepsilon)
=
1-0+1
=
2.
\]
This proves the assertion.
\end{proof}




\begin{proposition}\label{prop:real-part-link}
Let $f$ be a nonconstant holomorphic function defined near $0\in \mathbb C^2$, with
\[
        f(0)=0.
\]
Let
\[
        X=\{z\in\mathbb C^2:\operatorname{Re} f(z)=0\}
\]
as a germ at $0$. If the embedded germ $(\mathbb R^4,X,0)$ is homeomorphic, as a germ of pairs, to $(\mathbb R^4,H,0)$, where $H\subset\mathbb R^4$ is a real hyperplane, then
\[
        df(0)\ne 0.
\]
Equivalently, if $df(0)=0$, then the germ $(\mathbb R^4,X,0)$ is not locally homeomorphic, as an embedded pair, to a flat real hyperplane germ.
\end{proposition}


\begin{proof}
We prove the contrapositive. Assume that
\[
        df(0)=0
\]
and suppose, toward a contradiction, that \(X\) is locally flat at
\(0\).

After shrinking the representative neighborhood, factor
\[
        f=u f_1^{m_1}\cdots f_r^{m_r},
\]
where \(u\) is a nowhere-zero holomorphic unit, the \(f_j\) are
pairwise nonassociated irreducible germs, and \(m_j\ge1\).

\smallskip
\noindent\emph{The nonreduced case.}
Suppose first that \(m_j\ge2\) for some \(j\). Choose a point
\[
        p\in\{f_j=0\},\qquad p\ne0,
\]
arbitrarily close to \(0\), such that \(p\) is a smooth point of the
branch \(\{f_j=0\}\) and lies on no other branch. Such points exist
because an irreducible plane curve germ is smooth away from the
origin after shrinking the representative, and distinct local
branches meet only at the origin after a further shrinking.

Near \(p\), choose holomorphic coordinates \((\eta,w)\), centered at
\(p\), in which
\[
        \{f_j=0\}=\{w=0\}.
\]
All nonvanishing factors can be absorbed into a holomorphic unit, and
after taking a local \(m_j\)-th root of that unit and changing the
\(w\)-coordinate, we obtain
\[
        f(\eta,w)=w^{m_j}.
\]
Hence, near \(p\),
\[
        X=\{\operatorname{Re}(w^{m_j})=0\}.
\]
Writing \(w=re^{i\theta}\), the zero set in the transverse
\(w\)-plane consists of the \(2m_j\) half-rays
\[
        \theta=\frac{\pi/2+\ell\pi}{m_j},
        \qquad
        \ell=0,\ldots,2m_j-1.
\]
Therefore the germ of the complement of \(X\) at \(p\) has
\(2m_j\) local connected components: locally the \(\eta\)-direction
is a product factor, while the \(w\)-plane is divided into exactly
\(2m_j\) sectors by the above rays. Since \(m_j\ge2\), there are at
least four local complementary components.

On the other hand, local flatness at \(0\) means that there are
neighborhoods \(U,V\) of \(0\) and a homeomorphism of pairs
\[
        (U,X\cap U)\longrightarrow(V,H\cap V),
\]
where \(H\) is a real hyperplane. Choosing \(p\) sufficiently close
to \(0\), we have \(p\in U\). Restricting this homeomorphism to a
small neighborhood of \(p\) shows that \(X\) is locally flat at
\(p\). But the germ of the complement of a real hyperplane has
exactly two local connected components. This contradicts
\(2m_j\ge4\).

Thus the nonreduced case is impossible.

\smallskip
\noindent\emph{The reduced case.}
We are left with
\[
        f=u f_1\cdots f_r.
\]
Thus the plane curve germ \(f^{-1}(0)\) is reduced. Since
\(df(0)=0\), the curve is singular at the origin. We record why the
critical point of \(f\) is isolated. If the critical set of \(f\) had a
positive-dimensional analytic component \(C\) through \(0\), then
\(df\) would vanish identically on the smooth part of \(C\), so \(f\)
would be constant on \(C\). Since \(f(0)=0\), this would imply
\(C\subset f^{-1}(0)\). Every point of \(C\) would then be a singular
point of the curve \(f^{-1}(0)\), contradicting the fact that a reduced
plane curve germ has zero-dimensional singular set. Hence, after
shrinking the representative, \(0\) is the only critical point of
\(f\). Equivalently, the Jacobian ideal
\[
        J(f)=\left(
        \frac{\partial f}{\partial z_1},
        \frac{\partial f}{\partial z_2}
        \right)
        \subset \mathbb C\{z_1,z_2\}
\]
has finite codimension. Thus the Milnor number
\[
        \mu(f)=
        \dim_{\mathbb C}
        \frac{\mathbb C\{z_1,z_2\}}{J(f)}
\]
is finite. See, for example, \cite[Secs.~4--7]{Milnor1968},
\cite[Ch.~3]{Dimca1992}, and \cite[Ch.~5]{Wall2004}.

Choose a sufficiently small admissible sphere
\(S^3_\varepsilon\), and set
\[
        K=f^{-1}(0)\cap S^3_\varepsilon,
        \qquad
        L_\varepsilon=X\cap S^3_\varepsilon.
\]
Since \(X\) is assumed locally flat at \(0\),
Lemma~\ref{lem:flat-link-homology} gives
\[
        \chi(L_\varepsilon)=2.
        \tag{3.3}
\]

Milnor's fibration theorem gives
\[
        \frac{f}{|f|}:S^3_\varepsilon\setminus K\longrightarrow S^1.
\]
If \(P_\theta\) denotes the closure of the fiber with argument
\(\theta\), then
\[
        \partial P_\theta=K,
\]
and all pages are homeomorphic compact connected orientable
surfaces. Moreover, away from \(K\), the equation
\(\operatorname{Re}f=0\) is equivalent to
\[
        \arg f=\frac{\pi}{2}
        \quad\text{or}\quad
        \arg f=\frac{3\pi}{2}.
\]
Consequently,
\[
        L_\varepsilon
        =
        P_{\pi/2}\cup_K P_{3\pi/2}.
        \tag{3.4}
\]
Here \(P_{\pi/2}\) and \(P_{3\pi/2}\) are the closed pages of the
Milnor open book; their interiors are disjoint, and their intersection
is precisely their common boundary \(K\). Since \(K\) is a finite
disjoint union of circles,
\[
        \chi(K)=0.
\]
Using additivity of Euler characteristic in (3.4) and writing
\(P=P_{\pi/2}\), we obtain
\[
        \chi(L_\varepsilon)=2\chi(P).
        \tag{3.5}
\]

The page \(P\) is homotopy equivalent to the Milnor fiber of \(f\),
which has the homotopy type of a bouquet of \(\mu(f)\) circles.
Therefore
\[
        \chi(P)=1-\mu(f).
        \tag{3.6}
\]
Combining (3.3)--(3.6) gives
\[
        2
        =
        \chi(L_\varepsilon)
        =
        2(1-\mu(f)),
\]
and hence
\[
        \mu(f)=0.
\]

This is impossible under the assumption \(df(0)=0\). Indeed,
\[
        \frac{\partial f}{\partial z_1}(0)
        =
        \frac{\partial f}{\partial z_2}(0)
        =0,
\]
so both partial derivatives belong to the maximal ideal
\[
        \mathfrak m=(z_1,z_2)
        \subset\mathbb C\{z_1,z_2\}.
\]
Hence
\[
        \left(
        \frac{\partial f}{\partial z_1},
        \frac{\partial f}{\partial z_2}
        \right)
        \subset\mathfrak m
        \ne\mathbb C\{z_1,z_2\},
\]
and therefore
\[
        \mu(f)\ge1,
\]
contrary to \(\mu(f)=0\).

Both the nonreduced and the reduced cases lead to contradictions.
Thus \(X=\{\operatorname{Re}f=0\}\) cannot be locally flat at a
critical point of \(f\). This proves the proposition.
\end{proof}



\begin{remark}\label{rem:why-link-not-components}
The proof uses invariants forced by local flatness, rather than merely
the number of sectors seen in one-dimensional complex slices. In the
nonreduced case, an exact local product normal form shows directly that
the ambient complement has more than two local components. In the
reduced case, local flatness determines the local homology and therefore
the homology, in particular the Euler characteristic, of a sufficiently
small admissible link. For instance,
\[
        \{\re(z_1^2+z_2^2)=0\}
\]
has a torus as its link on a small sphere, and hence Euler
characteristic \(0\), whereas local flatness would force Euler
characteristic \(2\).
\end{remark}


\section{Proof of the Lewy-type theorem}

We now prove Theorem~\ref{th:main}.


\begin{proof}[Proof of Theorem~\ref{th:main}]
We give all details of the local reduction. The assertion is local at
$p$. Hence we may replace $\Omega$ by a smaller neighborhood of $p$ at
any point of the proof.

First translate the source and the target. Thus we may assume that
\[
        p=0,
        \qquad
        F(0)=0.
\]
Translations preserve pluriharmonicity, local injectivity, and the value
of the real Jacobian determinant up to evaluation at the translated
point. Choose a ball
\[
        B=B(0,\rho),
        \qquad
        \overline B\subset \Omega,
\]
such that $F$ is one-to-one on $B$. We identify $\C^2$ with $\R^4$
and write
\[
        F=(u_1,u_2,u_3,u_4):B\subset\R^4\longrightarrow\R^4,
\]
where the $u_j$ are real-valued $C^2$ functions. Since $F$ is
pluriharmonic, each $u_j$ is pluriharmonic.

Suppose, toward a contradiction, that
\[
        J_F(0)=0.
\]
The Jacobian $J_F(0)$ is the determinant of the real linear map
\[
        dF(0):\R^4\longrightarrow\R^4.
\]
Thus $dF(0)$ is not invertible. Its image is a proper real linear
subspace of $\R^4$. Therefore the annihilator of $\operatorname{im}dF(0)$
in the dual space $(\R^4)^*$ is nontrivial. Hence there exists a
nonzero real linear functional
\[
        \lambda:\R^4\longrightarrow\R
\]
such that
\[
        \lambda(dF(0)v)=0
        \qquad
        \hbox{for every }v\in\R^4.
\]
Equivalently,
\[
        \lambda\circ dF(0)=0.
\]
Set
\[
        u=\lambda\circ F.
\]
Then $u$ is a real-valued $C^2$ function on $B$. By the chain rule,
\[
        du(0)=d(\lambda\circ F)(0)=\lambda\circ dF(0)=0.
\]
Also,
\[
        u(0)=\lambda(F(0))=0.
\]

The function $u$ is pluriharmonic. Indeed, since $\lambda$ is real
linear, there are real constants $a_1,a_2,a_3,a_4$ such that
\[
        \lambda(x_1,x_2,x_3,x_4)
        =a_1x_1+a_2x_2+a_3x_3+a_4x_4.
\]
Consequently
\[
        u=a_1u_1+a_2u_2+a_3u_3+a_4u_4.
\]
The class of real-valued pluriharmonic functions is closed under real
linear combinations, so $u$ is pluriharmonic.

We next show that $u$ is not identically zero in any neighborhood of
$0$. Since $F$ is continuous and one-to-one on the open set $B$,
invariance of domain implies that $F(B)$ is open in $\R^4$. Moreover,
$F:B\to F(B)$ is a homeomorphism: it is a continuous bijection onto its
image, and it is an open map by invariance of domain. If $u$ were
identically zero on some ball $B'\subset B$ centered at $0$, then
\[
        F(B')\subset H,
        \qquad
        H:=\{x\in\R^4:\lambda(x)=0\}.
\]
Here $H$ is a real hyperplane because $\lambda\ne0$. But $F(B')$ is
open in $\R^4$, while $H$ has empty interior in $\R^4$. This is
impossible. Hence $u$ is nonconstant.

After shrinking $B$ if necessary, we may assume that $B$ is simply
connected. Since $u$ is pluriharmonic, Lemma~\ref{lem:primitive} gives
a holomorphic function $f$ on $B$ such that
\[
        u=\re f.
\]
Because $u(0)=0$, we have $\re f(0)=0$. Thus $f(0)$ is purely
imaginary. Replacing $f$ by $f-f(0)$, which does not change its real
part, we may assume that
\[
        f(0)=0.
\]
The function $f$ is nonconstant; otherwise $u=\re f$ would be constant,
contrary to the preceding paragraph.

We now prove that $df(0)=0$. Let $v\in T_0\R^4$ be an arbitrary real
tangent vector. Since $u=\re f$,
\[
        du(0)[v]=d(\re f)(0)[v]=\re(df(0)[v]).
\]
But $du(0)=0$, so
\[
        \re(df(0)[v])=0
        \qquad
        \hbox{for every }v\in T_0\R^4.
\]
Since $f$ is holomorphic, the real differential $df(0)$ is complex
linear. Thus
\[
        df(0)[iv]=i\,df(0)[v].
\]
Applying the preceding identity to the real tangent vector $iv$ gives
\[
        0
        =\re(df(0)[iv])
        =\re(i\,df(0)[v])
        =-\im(df(0)[v]).
\]
Therefore both the real part and the imaginary part of $df(0)[v]$
vanish. Since $v$ was arbitrary,
\[
        df(0)=0.
\]

It remains to record exactly where local flatness enters. Let
\[
        H=\{x\in\R^4:\lambda(x)=0\}.
\]
Then $H$ is a real hyperplane through the origin. Since $F:B\to F(B)$
is a homeomorphism and $F(0)=0$, the restriction of $F$ gives a
homeomorphism of germs of pairs
\[
        (B,F^{-1}(H),0)
        \longrightarrow
        (F(B),H\cap F(B),0).
\]
The pair on the right is exactly the germ at $0$ of the real hyperplane
$H$ in $\R^4$. Hence the embedded germ $F^{-1}(H)$ is locally flat at
$0$.

On the other hand,
\[
        F^{-1}(H)
        =\{z\in B:\lambda(F(z))=0\}
        =\{z\in B:u(z)=0\}
        =\{z\in B:\re f(z)=0\}.
\]
Thus the real hypersurface germ $\{\re f=0\}$ is locally flat at
$0$.

We have obtained a nonconstant holomorphic function $f$ near
$0\in\C^2$ such that
\[
        f(0)=0,
        \qquad
        df(0)=0,
\]
and such that the embedded real hypersurface germ $\{\re f=0\}$ is
locally flat at $0$. This contradicts Proposition~\ref{prop:real-part-link}.
The contradiction proves that $J_F(0)\ne0$. Translating back, we obtain
\[
        J_F(p)\ne0.
\]
This proves the theorem.
\end{proof}


\begin{corollary}\label{cor:univalent}
Let $\Omega\subset\C^2$ be a domain, and let $F:\Omega\to\C^2$ be a
univalent $C^2$ pluriharmonic mapping. Then
\[
        J_F(z)\ne0,
        \qquad z\in\Omega.
\]
Consequently, $F$ is a local $C^1$ diffeomorphism at every point of
$\Omega$.
\end{corollary}

\begin{proof}
Apply Theorem~\ref{th:main} at each point. The final assertion follows
from the inverse function theorem.
\end{proof}
\begin{remark}\label{rem:scope-and-context}
Naser \cite{Naser1973} stated a Lewy-type assertion for pluriharmonic
mappings. The proof given above supplies a detailed local argument in
the case of mappings \(\C^2\to\C^2\), emphasizing the reduction to the
real hypersurface
\[
        \{\re f=0\}
\]
and the topology of plane curve singularities.

The proof is specific to \(\C^2\). The pluriharmonicity assumption makes
a real linear projection of \(F\) locally equal to the real part of a
holomorphic function. In two complex variables, the critical zero set of
such a function is governed by the topology of plane curve singularities:
the link lies in \(S^3\), and Milnor fibers are compact orientable
surfaces. In higher complex dimension, the corresponding links and
Milnor fibers have higher dimension, and the argument above does not by
itself prove a Lewy theorem.

Wood's example \cite{Wood1991} shows that the direct higher-dimensional
analogue of Lewy's planar theorem is false for harmonic homeomorphisms.
Thus the theorem above should be viewed not as a consequence of ordinary
harmonicity alone, but as a consequence of the additional complex
structure encoded in pluriharmonicity.
\end{remark}

\section{Quasiconformal harmonic local diffeomorphisms onto Dini-smooth domains}

We now prove the bi-Lipschitz application.  The point of the formulation
below is that the analytic part needed in this section is only local
nondegeneracy: once a harmonic quasiconformal homeomorphism is known to be a
local diffeomorphism, Dini smoothness of the image domain gives the global
bi-Lipschitz conclusion.  The Lewy theorem proved in this paper is then used
only to verify the local-diffeomorphism hypothesis for pluriharmonic mappings
in \(\C^2\).  In the planar case the same hypothesis follows from Lewy's
classical theorem.

Throughout this section, if \(D\subset\R^m\) is a domain, we write
\[
        \delta_D(x)=\operatorname{dist}(x,\partial D).
\]
If \(A:\R^m\to\R^m\) is linear, set
\[
        |A|=\sup_{|v|=1}|Av|,
        \qquad
        l(A)=\inf_{|v|=1}|Av|.
\]

\begin{definition}[Dini modulus and \(C^1\)-Dini boundary]
\label{def:dini-boundary}
A modulus of continuity is a continuous increasing function
\[
        \omega:[0,\omega_0)\to[0,\infty),
        \qquad
        \omega(0)=0.
\]
It is called a Dini modulus if
\[
        \int_0^{\omega_0}\frac{\omega(t)}{t}\,dt<\infty .
\]
A bounded domain \(\Omega\subset\R^m\) has \(C^1\)-Dini boundary if, in a
finite system of boundary charts, \(\partial\Omega\) is represented as the
graph of a \(C^1\) function whose gradient has a common Dini modulus of
continuity.
\end{definition}

\begin{lemma}[Regularized defining function]
\label{lem:regularized-defining-function-dini}
Let \(\Omega\subset\R^m\) be a bounded domain with \(C^1\)-Dini boundary.
Let
\[
        d(y)=\operatorname{dist}(y,\partial\Omega),
        \qquad y\in\Omega .
\]
Then there exist constants \(\varepsilon_0>0\), \(C_0\ge1\), a Dini
modulus \(\omega\), and a function
\[
        \rho\in C^\infty(\Omega), \qquad \rho>0,
\]
such that, whenever \(0<d(y)<\varepsilon_0\),
\[
        C_0^{-1}d(y)\le \rho(y)\le C_0d(y),
\]
\[
        C_0^{-1}\le |\nabla\rho(y)|\le C_0,
\]
and
\[
        |D^2\rho(y)|
        \le
        C_0\frac{\omega(\rho(y))}{\rho(y)}.
\]
Moreover,
\[
        \int_0^{\varepsilon_0}\frac{\omega(t)}{t}\,dt<\infty .
\]
\end{lemma}

\begin{proof}
This is the regularized distance theorem of Lieberman
\cite{Lieberman1985}.  A \(C^1\)-Dini domain admits, near its boundary, a
\(C^1\) defining function whose gradient has a Dini modulus of continuity
and is bounded away from zero on the boundary.  Lieberman's construction
regularizes this defining function inside \(\Omega\), gives comparability
with the distance to the boundary, and yields the displayed second-derivative
estimate.  The Dini condition is preserved after changing the argument of the
modulus by fixed multiplicative constants.
\end{proof}

\begin{lemma}[Boundary-distance estimate]
\label{lem:boundary-distance-dini}
Let \(\Omega\subset\R^m\) be a bounded domain with \(C^1\)-Dini boundary,
and let
\[
        f:B^m\to\Omega
\]
be a harmonic \(K\)-quasiconformal homeomorphism which is a local
\(C^1\)-diffeomorphism in \(B^m\). Then there are constants \(c,C>0\) such
that
\[
        c(1-|x|)
        \le
        \delta_\Omega(f(x))
        \le
        C(1-|x|),
        \qquad x\in B^m .
\]
\end{lemma}

\begin{proof}
Let \(\rho\) be the regularized defining function from
Lemma~\ref{lem:regularized-defining-function-dini}. Since \(\Omega\) is
bounded and \(f\) is a homeomorphism,
\[
        \delta_\Omega(f(x))\to0
        \qquad\text{as } |x|\to1.
\]
Indeed, otherwise a sequence \(|x_j|\to1\) would have a subsequence with
\(f(x_j)\) converging to a point of \(\Omega\), contradicting continuity of
\(f^{-1}\).  Hence, for some \(r_0\in(0,1)\), the function
\[
        v(x)=\rho(f(x))
\]
is defined in the collar
\[
        A_{r_0}=\{x\in B^m:r_0<|x|<1\}.
\]

Since the coordinate functions of \(f\) are harmonic, the chain rule gives
\[
        \Delta v(x)
        =
        \sum_{\alpha,\beta=1}^m
        \rho_{\alpha\beta}(f(x))
        \langle \nabla f^\alpha(x),\nabla f^\beta(x)\rangle .
\]
Therefore
\[
        |\Delta v(x)|
        \le
        C |D^2\rho(f(x))|\,|Df(x)|^2 .
\]
By Lemma~\ref{lem:regularized-defining-function-dini},
\[
        |D^2\rho(f(x))|
        \le
        C\frac{\omega(v(x))}{v(x)}.
\]
Since \(f\) is a local diffeomorphism and \(K\)-quasiconformal, its
differential satisfies
\[
        |Df(x)|\le K_1 l(Df(x)),
\]
where \(K_1\) depends only on \(K\) and \(m\).  Moreover,
\[
\begin{aligned}
        |\nabla v(x)|
        &=|Df(x)^T\nabla\rho(f(x))|       \\
        &\ge l(Df(x))|\nabla\rho(f(x))|  \\
        &\ge C^{-1}l(Df(x)).
\end{aligned}
\]
Consequently,
\[
        |\Delta v(x)|
        \le
        A\frac{\omega(v(x))}{v(x)}|\nabla v(x)|^2,
        \qquad x\in A_{r_0},
\]
for a constant \(A>0\).

Put
\[
        \sigma(t)=\frac{\omega(t)}{t},
        \qquad
        I(t)=\int_0^t \sigma(s)\,ds.
\]
The Dini condition gives \(I(t)<\infty\) for small \(t\) and
\(I(t)\to0\) as \(t\downarrow0\).  Define
\[
        \Phi_-(t)=\int_0^t e^{-AI(s)}\,ds,
        \qquad
        \Phi_+(t)=\int_0^t e^{AI(s)}\,ds .
\]
After decreasing \(\varepsilon_0\), if necessary, both \(\Phi_-\) and
\(\Phi_+\) are comparable to \(t\) on \((0,\varepsilon_0)\).  Also,
\[
        \Phi_-''(t)=-A\sigma(t)\Phi_-'(t),
        \qquad
        \Phi_+''(t)=A\sigma(t)\Phi_+'(t).
\]

Set
\[
        w_-(x)=\Phi_-(v(x)),
        \qquad
        w_+(x)=\Phi_+(v(x)).
\]
Using the differential inequality for \(v\), we get
\[
        \Delta w_-\le0,
        \qquad
        \Delta w_+\ge0
\]
in \(A_{r_0}\).  Thus \(w_-\) is positive superharmonic and \(w_+\) is
positive subharmonic.


Let \(h\) be the radial harmonic function in
\[
        A_{r_0}=\{x\in \mathbb R^m:r_0<|x|<1\}
\]
which equals \(1\) on \(|x|=r_0\) and \(0\) on \(|x|=1\).  Writing
\(r=|x|\), this function is given explicitly by
\[
h(x)=H(r),
\]
where
\[
H(r)=
\begin{cases}
\dfrac{\log r}{\log r_0}
      =\dfrac{\log(1/r)}{\log(1/r_0)}, & m=2, \\[1.2em]
\dfrac{r^{2-m}-1}{r_0^{2-m}-1}, & m\ge 3 .
\end{cases}
\]
Indeed, these are precisely the radial harmonic functions in dimensions
\(2\) and \(m\ge 3\), respectively, and the constants are chosen so that
\(H(r_0)=1\) and \(H(1)=0\).

We shall use that
\[
        h(x)\asymp 1-|x|,
        \qquad r_0<|x|<1 .
\]
For \(m=2\), this follows from
\[
        \log(1/r)=\int_r^1 \frac{dt}{t},
\]
and hence, since \(r_0<r<1\),
\[
        1-r \le \log(1/r)\le \frac{1}{r_0}(1-r).
\]
Therefore
\[
        \frac{1}{\log(1/r_0)}(1-r)
        \le H(r)
        \le
        \frac{1}{r_0\log(1/r_0)}(1-r).
\]
For \(m\ge 3\), we have
\[
        r^{2-m}-1
        =
        \int_r^1 (m-2)t^{1-m}\,dt .
\]
Since \(r_0<r<t<1\), this gives
\[
        (m-2)(1-r)
        \le r^{2-m}-1
        \le
        (m-2)r_0^{1-m}(1-r).
\]
Consequently,
\[
        \frac{m-2}{r_0^{2-m}-1}(1-r)
        \le H(r)
        \le
        \frac{(m-2)r_0^{1-m}}{r_0^{2-m}-1}(1-r).
\]
Thus, in all dimensions \(m\ge 2\),
\[
        h(x)\asymp 1-|x|,
        \qquad r_0<|x|<1,
\]
with constants depending only on \(m\) and \(r_0\).


The functions \(w_-\) and \(w_+\) extend continuously by zero to the
outer boundary. On the compact inner sphere \(|x|=r_0\), set
\[
        m_-:=\min_{|x|=r_0}w_-(x)>0,
        \qquad
        M_+:=\max_{|x|=r_0}w_+(x)<\infty.
\]
Since \(h=1\) on \(|x|=r_0\) and \(h=0\) on \(|x|=1\), the minimum
principle applied to \(w_- -m_-h\) and the maximum principle applied to
\(w_+-M_+h\) give
\[
        m_-h(x)\le w_-(x),
        \qquad
        w_+(x)\le M_+h(x).
\]
Because \(\Phi_\pm(t)\asymp t\) and \(h(x)\asymp1-|x|\), it follows
that
\[
        c(1-|x|)\le v(x)\le C(1-|x|),
        \qquad r_0<|x|<1 .
\]
Since \(v=\rho\circ f\) and \(\rho\asymp\delta_\Omega\), this proves the
claimed estimate near the boundary.  On the compact ball
\(\overline{B^m_{r_0}}\) the estimate follows after changing the constants.
\end{proof}

\begin{lemma}[Lipschitz estimate]
\label{lem:lipschitz-dini}
Let \(\Omega\subset\R^m\) be a bounded \(C^1\)-Dini domain, and let
\[
        f:B^m\to\Omega
\]
be a harmonic \(K\)-quasiconformal homeomorphism which is a local
\(C^1\)-diffeomorphism. Then \(f\) is Lipschitz.
\end{lemma}

\begin{proof}
By Lemma~\ref{lem:boundary-distance-dini},
\[
        \delta_\Omega(f(x))\le C(1-|x|).
\]
Fix \(x\in B^m\), and put
\[
        r=\frac{1-|x|}{2}.
\]
If \(y\in B(x,r)\), then \(B(x,r)\subset B^m\) and
\[
        k_{B^m}(x,y)\le C.
\]
By the quasi-invariance of the quasihyperbolic metric under
quasiconformal mappings, due to Gehring and Osgood
\cite{GehringOsgood1979},
\[
        k_\Omega(f(x),f(y))\le C.
\]
For every domain \(D\) one has
\[
        k_D(a,b)
        \ge
        \log\left(1+\frac{|a-b|}{\delta_D(a)}\right).
\]
Therefore
\[
        |f(y)-f(x)|
        \le
        C\delta_\Omega(f(x))
        \le
        C(1-|x|),
        \qquad y\in B(x,r).
\]
The interior gradient estimate for harmonic functions gives
\[
        |Df(x)|
        \le
        \frac{C}{r}\sup_{B(x,r)}|f(y)-f(x)|
        \le C.
\]
Thus \(|Df|\) is bounded in \(B^m\), and \(f\) is Lipschitz.
\end{proof}

\begin{lemma}[A half-space Liouville lemma]
\label{lem:halfspace-liouville}
Let
\[
        H=\{x=(x',x_m)\in\R^m:x_m>0\}.
\]
If \(u\) is a positive harmonic function in \(H\) and
\[
        0<u(x)\le Cx_m,
        \qquad x\in H,
\]
then
\[
        u(x)=a x_m
\]
for some constant \(a\ge0\).  If in addition \(u(e_m)>0\), then
\(a>0\).
\end{lemma}

\begin{proof}
The estimate \(u(x)\le Cx_m\) implies that \(u\) extends continuously to
\(\partial H\) by the value \(0\).  By the Schwarz reflection principle for
harmonic functions, the odd extension
\[
        U(x',x_m)=
        \begin{cases}
        u(x',x_m), & x_m>0,\\
        -u(x',-x_m), & x_m<0,\\
        0, & x_m=0,
        \end{cases}
\]
is harmonic in all of \(\R^m\).  Moreover,
\[
        |U(x)|\le C|x_m|\le C|x|.
\]
Thus \(U\) is an entire harmonic function of at most linear growth.  By
the standard Liouville theorem for harmonic functions of polynomial growth,
\(U\) is affine.  Since \(U=0\) on \(\{x_m=0\}\) and \(U\) is odd in
\(x_m\), it follows that \(U(x)=a x_m\).  Hence \(u(x)=a x_m\) in \(H\).
If \(u(e_m)>0\), then \(a=u(e_m)>0\).
\end{proof}

\begin{lemma}[Quasiconvexity of \(C^1\) domains]
\label{lem:c1-quasiconvex}
Let \(\Omega\subset\R^m\) be a bounded connected domain with \(C^1\)
boundary. Then there is a constant \(M_\Omega\ge1\) such that any two
points \(p,q\in\Omega\) can be joined by a rectifiable curve
\(\gamma\subset\Omega\) satisfying
\[
        \ell(\gamma)\le M_\Omega |p-q|.
\]
\end{lemma}

\begin{proof}
Bounded \(C^1\) domains are Lipschitz domains, and bounded connected
Lipschitz domains are quasiconvex. See, for example, Martio and Sarvas
\cite{MartioSarvas1979}.
\end{proof}

\begin{theorem}[Bi-Lipschitz harmonic quasiconformal local diffeomorphisms]
\label{th:harmonic-qc-local-diffeo-bilip}
Let \(\Omega\subset\R^m\) be a bounded domain with \(C^1\)-Dini boundary,
and let
\[
        f:B^m\to\Omega
\]
be a harmonic \(K\)-quasiconformal homeomorphism. Assume that \(f\) is a
local \(C^1\)-diffeomorphism in \(B^m\). Then \(f\) is bi-Lipschitz.
\end{theorem}

\begin{proof}
By Lemma~\ref{lem:lipschitz-dini}, \(f\) is Lipschitz.  Hence there is
\(L<\infty\) such that
\[
        |f(x)-f(y)|\le L|x-y|,
        \qquad x,y\in B^m.
\]

We prove that the minimal stretch has a positive lower bound. Suppose, to
the contrary, that there are points \(x_k\in B^m\) such that
\[
        l(Df(x_k))\to0.
\]
Since \(f\) is a local diffeomorphism, \(l(Df)>0\) in \(B^m\).  By
continuity, \(l(Df)\) has a positive minimum on compact subsets.  Hence
\[
        |x_k|\to1.
\]

Let
\[
        \rho_k=1-|x_k|,
        \qquad
        \xi_k=\frac{x_k}{|x_k|}.
\]
After passing to a subsequence, \(\xi_k\to\xi_0\in\partial B^m\).  Let
\(q_k=\xi_k\in\partial B^m\), and choose orthogonal maps \(U_k\) such that
\[
        U_k e_m=-\xi_k.
\]
Thus
\[
        x_k=q_k+\rho_k U_k e_m.
\]
Define the rescaled source domains
\[
        D_k=\{\zeta\in\R^m:q_k+\rho_k U_k\zeta\in B^m\}.
\]
Then \(D_k\) converges on compact subsets to the half-space
\[
        H=\{\zeta_m>0\},
\]
and the point \(e_m\in H\) corresponds to \(x_k\).

Let \(p_k=f(x_k)\), and choose \(a_k\in\partial\Omega\) with
\[
        |p_k-a_k|=\delta_\Omega(p_k).
\]
After passing to a subsequence, the boundary points \(a_k\) converge to
some \(a_0\in\partial\Omega\).  Since \(\partial\Omega\) is \(C^1\)-Dini,
there are orthogonal maps \(R_k\) sending the inward unit normal to
\(\partial\Omega\) at \(a_k\) to \(e_m\), such that the rescaled target
domains
\[
        \Omega_k=\frac{R_k(\Omega-a_k)}{\rho_k}
\]
converge on compact subsets to the half-space \(H\).  Define
\[
        f_k(\zeta)=
        \frac{R_k(f(q_k+\rho_kU_k\zeta)-a_k)}{\rho_k},
        \qquad \zeta\in D_k.
\]
At a nearest boundary point of a \(C^1\) boundary one has
\[
        p_k-a_k=\delta_\Omega(p_k)n(a_k),
\]
where \(n(a_k)\) is the inward unit normal. Hence, by the choice of
\(R_k\),
\[
        f_k(e_m)=\frac{\delta_\Omega(p_k)}{\rho_k}e_m.
\]
Lemma~\ref{lem:boundary-distance-dini} shows that these base values stay
in a fixed compact subinterval of the positive \(e_m\)-axis. Each
\(f_k\) is harmonic and \(K\)-quasiconformal. Together with the
Lipschitz bound for \(f\), the basepoint normalization implies that the
maps \(f_k\) are locally uniformly bounded and locally uniformly
Lipschitz on compact subsets of \(H\).
Harmonic interior estimates and Arzel\`a--Ascoli therefore give a
subsequence converging in \(C^1\) on compact subsets of \(H\) to a
harmonic mapping
\[
        f_\infty:H\to \overline H.
\]
Moreover,
\[
        l(Df_\infty(e_m))=
        \lim_{k\to\infty} l(Df_k(e_m))
        =
        \lim_{k\to\infty} l(Df(x_k))=0.
\]

Let
\[
        u(\zeta)=(f_\infty(\zeta))_m
\]
be the last coordinate of \(f_\infty\).  The two-sided estimate of
Lemma~\ref{lem:boundary-distance-dini}, after rescaling, reads
\[
        c\,\delta_{D_k}(\zeta)
        \le \delta_{\Omega_k}(f_k(\zeta))
        \le C\,\delta_{D_k}(\zeta).
\]
On compact subsets of the limiting half-space,
\(\delta_{D_k}(\zeta)\to\zeta_m\), while the \(C^1\) convergence of the
rescaled target boundaries gives
\(\delta_{\Omega_k}(w)\to w_m\) locally uniformly. Passing to the limit
therefore gives
\[
        c\zeta_m\le u(\zeta)\le C\zeta_m,
        \qquad \zeta\in H.
\]
In particular, \(u\) is a positive harmonic function in \(H\).  By
Lemma~\ref{lem:halfspace-liouville},
\[
        u(\zeta)=a\zeta_m
\]
with \(a>0\).  Hence
\[
        |Df_\infty(e_m)|\ge |\nabla u(e_m)|=a>0.
\]
Thus \(f_\infty\) is nonconstant. By the compactness theorem for
quasiconformal mappings, a nonconstant locally uniform limit of
\(K\)-quasiconformal homeomorphisms is again \(K\)-quasiconformal, indeed a
homeomorphism onto its image; see V\"ais\"al\"a \cite{Vaisala1971}. Since
\(f_\infty\) is \(C^1\), the pointwise distortion inequality gives
\[
        |Df_\infty(e_m)|\le K_1 l(Df_\infty(e_m)),
\]
with \(K_1\) depending only on \(K\) and \(m\).  Therefore
\[
        l(Df_\infty(e_m))\ge a/K_1>0,
\]
contradicting \(l(Df_\infty(e_m))=0\).  Thus
\[
        \inf_{x\in B^m}l(Df(x))>0.
\]

Choose \(c_1>0\) such that
\[
        l(Df(x))\ge c_1,
        \qquad x\in B^m.
\]
Then \(f^{-1}\) is locally \(C^1\) in \(\Omega\), and
\[
        |D(f^{-1})(f(x))|\le c_1^{-1}.
\]
Let \(x,y\in B^m\), and put \(p=f(x)\), \(q=f(y)\).  By
Lemma~\ref{lem:c1-quasiconvex}, there is a rectifiable curve
\(\gamma\subset\Omega\) joining \(p\) to \(q\) such that
\[
        \ell(\gamma)\le M_\Omega |p-q|.
\]
Therefore
\[
\begin{aligned}
        |x-y|
        &=|f^{-1}(p)-f^{-1}(q)|                         \\
        &\le \int_\gamma |D(f^{-1})(\eta)|\,ds(\eta)       \\
        &\le c_1^{-1}\ell(\gamma)                         \\
        &\le c_1^{-1}M_\Omega |p-q|.
\end{aligned}
\]
Thus
\[
        |f(x)-f(y)|\ge \frac{c_1}{M_\Omega}|x-y|.
\]
Together with the Lipschitz estimate, this proves that \(f\) is
bi-Lipschitz.
\end{proof}


\begin{corollary}[Small-distortion harmonic quasiconformal mappings]
\label{cor:small-distortion-dini-bilip}
Let \(\Omega\subset\R^m\) be a bounded domain with \(C^1\)-Dini
boundary, and let
\[
        f:B^m\to\Omega
\]
be a harmonic quasiconformal homeomorphism. Assume that its outer
dilatation satisfies
\[
        K_O(f)<3^{m-1}.
\]
Then \(f\) is bi-Lipschitz.
\end{corollary}

\begin{proof}
By the theorem of Bo\v zin and Mateljevi\'c
\cite{BozinMateljevic2015}, the small-distortion condition
\(K_O(f)<3^{m-1}\) excludes interior zeros of the Jacobian of a harmonic
quasiconformal mapping. Hence
\[
        J_f(x)\neq 0,
        \qquad x\in B^m .
\]
Since the coordinate functions of \(f\) are harmonic, \(f\) is smooth in
\(B^m\). Therefore the inverse function theorem implies that \(f\) is a
local \(C^1\)-diffeomorphism in \(B^m\). The conclusion now follows
from Theorem~\ref{th:harmonic-qc-local-diffeo-bilip}.
\end{proof}

\begin{corollary}[Quasiconformal pluriharmonic mappings in \(\C^2\)]
\label{cor:pluriharmonic-c2-bilip}
Let \(\Omega\subset\C^2\cong\R^4\) be a bounded domain with \(C^1\)-Dini
boundary, and let
\[
        F:\mathbb B^2\to\Omega
\]
be a \(K\)-quasiconformal pluriharmonic homeomorphism. Then \(F\) is
bi-Lipschitz. In particular, every quasiconformal pluriharmonic
homeomorphism of \(\mathbb B^2\) onto itself is bi-Lipschitz.
\end{corollary}

\begin{proof}
By Theorem~\ref{th:main}, the real Jacobian of \(F\) does not vanish in
\(\mathbb B^2\), because \(F\) is locally one-to-one.  Hence \(F\) is a
local \(C^1\)-diffeomorphism.  Since the real coordinate functions of a
pluriharmonic mapping are harmonic, Theorem~\ref{th:harmonic-qc-local-diffeo-bilip}
with \(m=4\) applies.
\end{proof}

\subsection{The planar case}

In the plane, the local-diffeomorphism hypothesis follows from Lewy's
classical theorem.  We record the result in a form where both the source and
the image are allowed to be finitely connected \(C^1\)-Dini domains.

\begin{corollary}[Planar Dini-smooth source and image]
\label{cor:planar-dini-bilip}
Let \(D,\Omega\subset\C\) be bounded finitely connected domains whose
boundary components are \(C^1\)-Dini Jordan curves. Let
\[
        f:D\to\Omega
\]
be a harmonic \(K\)-quasiconformal homeomorphism. Then \(f\) is
bi-Lipschitz.
\end{corollary}

\begin{proof}
By Koebe's circle-domain theorem there is a circular domain \(G\) and a
conformal map
\[
        \varphi:G\to D.
\]
Since \(\partial D\) is \(C^1\)-Dini, the Kellogg--Warschawski theorem
implies that \(\varphi\) extends to a \(C^1\)-diffeomorphism of the
closures and satisfies
\[
        0<c\le |\varphi'|\le C<\infty
        \qquad\text{on }\overline G .
\]
Thus \(\varphi\) is bi-Lipschitz.  It is enough to prove that
\[
        F=f\circ\varphi:G\to\Omega
\]
is bi-Lipschitz.  Since harmonicity and quasiconformality are preserved
under conformal changes of variables in the plane, \(F\) is harmonic and
\(K\)-quasiconformal.

By Lewy's classical theorem, \(F\) has nonvanishing Jacobian in \(G\), and
therefore \(F\) is a local \(C^1\)-diffeomorphism.

We verify the ingredients of
Theorem~\ref{th:harmonic-qc-local-diffeo-bilip} for the circular source
\(G\). Since \(G\) has only finitely many boundary components and each of
them is a circle, every boundary component admits a collar in which a
radial harmonic barrier is comparable to \(\delta_G\). The regularized
defining function of \(\Omega\), composed with \(F\), satisfies
\[
        |\Delta v|\le A\frac{\omega(v)}{v}|\nabla v|^2,
\]
and the same Dini changes of variables as above produce a superharmonic
and a subharmonic comparison function. Comparing them with the annular
barriers on each boundary collar gives
\[
        c\delta_G(z)\le \delta_\Omega(F(z))\le C\delta_G(z),
        \qquad z\in G.
\]
On the compact part of \(G\) away from the boundary the same estimate
holds after changing the constants. The quasihyperbolic argument and the
interior gradient estimate then give \(|DF|\le C\) throughout \(G\).

It remains to exclude degeneration of the minimal stretch near the
boundary. Suppose that \(z_k\to\partial G\) and
\(l(DF(z_k))\to0\), and put \(\rho_k=\delta_G(z_k)\). Because the
boundary of \(G\) is circular, after translating, rotating, and scaling
by \(\rho_k^{-1}\), the source domains converge on compact subsets to a
half-plane. The \(C^1\)-Dini regularity of \(\partial\Omega\) gives the
same local convergence of the rescaled target domains to a half-plane.
Normalize at nearest boundary points exactly as in the proof of
Theorem~\ref{th:harmonic-qc-local-diffeo-bilip}. The rescaled maps are
locally uniformly bounded and locally uniformly Lipschitz, so harmonic
interior estimates give a subsequential \(C^1_{\rm loc}\)-limit. At this
stage the limit is only known to be harmonic. The rescaled two-sided
boundary-distance estimate passes to the limit and shows that its height
coordinate is bounded above and below by positive constant multiples of
the height variable. By Lemma~\ref{lem:halfspace-liouville}, that height
coordinate is a nonzero linear function. Hence the limiting map is
nonconstant, and the compactness theorem for quasiconformal mappings now
implies that it is \(K\)-quasiconformal. The pointwise distortion
inequality then gives a positive lower bound for its minimal stretch at
the base point, contradicting \(l(DF(z_k))\to0\). Thus
\(\inf_G l(DF)>0\).

Thus \(F\) is bi-Lipschitz. Since \(\varphi\) is bi-Lipschitz, the original
map \(f=F\circ\varphi^{-1}\) is bi-Lipschitz as well.
\end{proof}

\begin{remark}\label{rem:dini-and-local-diffeomorphism}
The proof separates the roles of boundary regularity and local
nondegeneracy.  The \(C^1\)-Dini smoothness of the image domain gives the
regularized defining function and the two-sided boundary-distance estimate.
The assumption that the harmonic quasiconformal map is a local
diffeomorphism gives pointwise nondegeneracy in the interior.  The blow-up
argument and the half-space Liouville lemma then prevent this
nondegeneracy from degenerating at the boundary.

For pluriharmonic mappings in \(\C^2\), the local-diffeomorphism hypothesis
is a consequence of Theorem~\ref{th:main}.  In the plane, it is a consequence
of Lewy's classical theorem.  In higher real dimensions it is not a
consequence of ordinary harmonicity alone, as shown by Wood's examples.
\end{remark}


\begin{remark}[Historical context for harmonic quasiconformal mappings]
\label{rem:historical-hqc}
The metric regularity of harmonic quasiconformal mappings has a substantial
history. Martio's work on harmonic quasiconformal mappings of the unit disk
is one of the early systematic sources of the subject; in particular it
connected quasiconformality of the Poisson extension with quantitative
boundary regularity conditions \cite{Martio1969}. A decisive planar result
was later obtained by Pavlovi\'c, who proved that a harmonic homeomorphism of
the unit disk onto itself is quasiconformal if and only if it is
bi-Lipschitz in the Euclidean metric \cite{Pavlovic2002}.

Subsequent work considered more general target domains and boundary
regularity assumptions. Kalaj proved Lipschitz and related boundary-distance
estimates for harmonic quasiconformal mappings between Jordan domains and,
in particular, for Dini-smooth targets \cite{Kalaj2008,Kalaj2015Dini};
related boundary regularity phenomena, including Muckenhoupt-weight
properties and a Lindel\"of-type theorem for harmonic quasiconformal
mappings, were studied in \cite{Kalaj2015AIM}. The sharpness of
Dini-type assumptions is also reflected by examples in the merely
\(C^1\)-smooth case \cite{Kalaj2022RMI}.

 Astala and Manojlovi\'c
studied higher-dimensional counterparts of Pavlovi\'c's theorem and proved
Lipschitz estimates for harmonic quasiconformal maps in space under suitable
additional hypotheses \cite{AstalaManojlovic2015}.

On the lower Lipschitz side, Bo\v zin and Mateljevi\'c proved
co-Lipschitz estimates for harmonic quasiconformal mappings of the disk onto
Lyapunov Jordan domains \cite{BozinMateljevic2020}. In higher dimensions
they also proved a Jacobian nondegeneracy theorem for harmonic
quasiconformal mappings under the small-distortion condition
\[
        K_O<3^{m-1}.
\]
More precisely, their normal-family argument shows that if the Jacobian
vanishes at an interior point, then the first nonzero homogeneous term in
the Taylor expansion has degree at least \(3\); comparison with the
quasiconformal distortion estimate then forces \(K_O\ge 3^{m-1}\)
\cite{BozinMateljevic2015}. Consequently, their result supplies the
local-diffeomorphism hypothesis in
Theorem~\ref{th:harmonic-qc-local-diffeo-bilip} whenever
\(K_O<3^{m-1}\), and yields
Corollary~\ref{cor:small-distortion-dini-bilip} for harmonic
quasiconformal homeomorphisms of the unit ball onto bounded \(C^1\)-Dini
domains.

The result above fits into this line of work. Its boundary part follows the
Dini-regularized-distance and quasihyperbolic-metric method, while the lower
Lipschitz bound is obtained by a blow-up argument. The new point in the
present paper is that, in complex dimension two, the needed interior
nondegeneracy for pluriharmonic maps is supplied by the Lewy-type theorem
proved in Theorem~\ref{th:main}. Thus, for quasiconformal pluriharmonic
homeomorphisms in \(\C^2\), no smallness assumption on \(K\) is needed. In
the planar corollary the same role is played by Lewy's classical theorem.
\end{remark}
\subsection*{Funding}
The author is partially supported by the Ministry of Education, Science and Innovation of Montenegro through the grants \emph{Mathematical Analysis, Optimization and Machine Learning} and \emph{Complex-analytic and geometric techniques for non-Euclidean machine learning: theory and applications}.

\begin{thebibliography}{99}

\bibitem{AstalaManojlovic2015}
K. Astala and V. Manojlovi\'c,
\newblock On Pavlovi\'c's theorem in space,
\newblock \emph{Potential Analysis} \textbf{43} (2015), no. 3, 361--370.
\newblock DOI: \texttt{10.1007/s11118-015-9475-4}.


\bibitem{BochnakCosteRoy1998}
J. Bochnak, M. Coste and M.-F. Roy,
\newblock \emph{Real Algebraic Geometry},
\newblock Ergebnisse der Mathematik und ihrer Grenzgebiete, 3. Folge, Vol. 36,
\newblock Springer-Verlag, Berlin, 1998.


\bibitem{BozinMateljevic2015}
V. Bo\v zin and M. Mateljevi\'c,
\newblock Bounds for Jacobian of harmonic injective mappings in
\(n\)-dimensional space,
\newblock \emph{Filomat} \textbf{29} (2015), no. 9, 2119--2124.
\newblock DOI: \texttt{10.2298/FIL1509119B}.

\bibitem{BozinMateljevic2020}
V. Bo\v zin and M. Mateljevi\'c,
\newblock Quasiconformal and HQC mappings between Lyapunov Jordan domains,
\newblock \emph{Annali della Scuola Normale Superiore di Pisa, Classe di
Scienze} (5) \textbf{21} (2020), 107--132.
\newblock DOI: \texttt{10.2422/2036-2145.201708\_013}.

\bibitem{BrieskornKnorrer1986}
E. Brieskorn and H. Kn\"orrer,
\newblock \emph{Plane Algebraic Curves},
\newblock translated by J. Stillwell,
\newblock Birkh\"auser, Basel, 1986.

\bibitem{BurgheleaVerona1972}
D. Burghelea and A. Verona,
\newblock Local homological properties of analytic sets,
\newblock \emph{Manuscripta Mathematica} \textbf{7} (1972), 55--66.

\bibitem{Chirka1989}
E. M. Chirka,
\newblock \emph{Complex Analytic Sets},
\newblock Kluwer Academic Publishers, Dordrecht, 1989.

\bibitem{Dimca1992}
A. Dimca,
\newblock \emph{Singularities and Topology of Hypersurfaces},
\newblock Universitext,
\newblock Springer-Verlag, New York, 1992.

\bibitem{GehringOsgood1979}
F. W. Gehring and B. G. Osgood,
\newblock Uniform domains and the quasi-hyperbolic metric,
\newblock \emph{Journal d'Analyse Math\'ematique} \textbf{36} (1979),
50--74.

\bibitem{GrauertRemmert1984}
H. Grauert and R. Remmert,
\newblock \emph{Coherent Analytic Sheaves},
\newblock Grundlehren der mathematischen Wissenschaften, Vol. 265,
\newblock Springer-Verlag, Berlin, 1984.

\bibitem{Hormander1990}
L. H\"ormander,
\newblock \emph{An Introduction to Complex Analysis in Several Variables},
\newblock third revised edition,
\newblock North-Holland, Amsterdam, 1990.


\bibitem{Kalaj2008}
D. Kalaj,
\newblock Quasiconformal harmonic mappings between Jordan domains,
\newblock \emph{Mathematische Zeitschrift} \textbf{260} (2008), 237--252.

\bibitem{Kalaj2015Dini}
D. Kalaj,
\newblock Quasiconformal harmonic mappings between Dini-smooth Jordan domains,
\newblock \emph{Pacific Journal of Mathematics} \textbf{276} (2015), no. 1,
213--228.
\newblock DOI: \texttt{10.2140/pjm.2015.276.213}.

\bibitem{Kalaj2015AIM}
D. Kalaj,
\newblock Muckenhoupt weights and Lindel\"of theorem for harmonic mappings,
\newblock \emph{Advances in Mathematics} \textbf{280} (2015), 301--321.
\newblock DOI: \texttt{10.1016/j.aim.2015.04.025}.



\bibitem{Kalaj2022RMI}
D. Kalaj,
\newblock Harmonic quasiconformal mappings between \(C^1\) smooth Jordan domains,
\newblock \emph{Revista Matem\'atica Iberoamericana} \textbf{38} (2022), no. 1,
113--130.
\newblock DOI: \texttt{10.4171/RMI/1272}.

\bibitem{Lewy1936}
H. Lewy,
\newblock On the non-vanishing of the Jacobian in certain one-to-one mappings,
\newblock \emph{Bulletin of the American Mathematical Society} \textbf{42} (1936), no. 10, 689--692.
\newblock DOI: \texttt{10.1090/S0002-9904-1936-06397-4}.

\bibitem{Lewy1968}
H. Lewy,
\newblock On the non-vanishing of the Jacobian of a homeomorphism by harmonic gradients,
\newblock \emph{Annals of Mathematics} \textbf{88} (1968), no. 3, 518--529.
\newblock DOI: \texttt{10.2307/1970723}.

\bibitem{Lieberman1985}
G. M. Lieberman,
\newblock Regularized distance and its applications,
\newblock \emph{Pacific Journal of Mathematics} \textbf{117} (1985), no. 2,
329--352.

\bibitem{Lojasiewicz1964}
S. \L{}ojasiewicz,
\newblock Triangulation of semi-analytic sets,
\newblock \emph{Annali della Scuola Normale Superiore di Pisa, Classe di Scienze} (3) \textbf{18} (1964), 449--474.


\bibitem{Martio1969}
O. Martio,
\newblock On harmonic quasiconformal mappings,
\newblock \emph{Annales Academiae Scientiarum Fennicae. Series A I. Mathematica}
No. 425 (1969), 1--10.
\newblock DOI: \texttt{10.5186/aasfm.1969.425}.

\bibitem{MartioSarvas1979}
O. Martio and J. Sarvas,
\newblock Injectivity theorems in plane and space,
\newblock \emph{Annales Academiae Scientiarum Fennicae. Series A I.
Mathematica} \textbf{4} (1979), 383--401.

\bibitem{Matsumura1986}
H. Matsumura,
\newblock \emph{Commutative Ring Theory},
\newblock Cambridge Studies in Advanced Mathematics, Vol. 8,
\newblock Cambridge University Press, Cambridge, 1986.

\bibitem{Milnor1968}
J. Milnor,
\newblock \emph{Singular Points of Complex Hypersurfaces},
\newblock Annals of Mathematics Studies, No. 61,
\newblock Princeton University Press, Princeton, 1968.

\bibitem{Naser1973}
M. Naser,
\newblock A generalization of a theorem of Lewy on harmonic mappings,
\newblock \emph{Mathematical Notes of the Academy of Sciences of the USSR} \textbf{13} (1973), no. 2, 113--115.
\newblock DOI: \texttt{10.1007/BF01094227}.

\bibitem{Pavlovic2002}
M. Pavlovi\'c,
\newblock Boundary correspondence under harmonic quasiconformal homeomorphisms
of the unit disk,
\newblock \emph{Annales Academiae Scientiarum Fennicae. Mathematica}
\textbf{27} (2002), 365--372.

\bibitem{Pommerenke1992}
Ch. Pommerenke,
\newblock \emph{Boundary Behaviour of Conformal Maps},
\newblock Grundlehren der mathematischen Wissenschaften, Vol. 299,
\newblock Springer-Verlag, Berlin, 1992.
\bibitem{Vaisala1971}
J. V\"ais\"al\"a,
\newblock \emph{Lectures on \(n\)-Dimensional Quasiconformal Mappings},
\newblock Lecture Notes in Mathematics, Vol. 229,
\newblock Springer-Verlag, Berlin--New York, 1971.

\bibitem{Wall2004}
C. T. C. Wall,
\newblock \emph{Singular Points of Plane Curves},
\newblock London Mathematical Society Student Texts, Vol. 63,
\newblock Cambridge University Press, Cambridge, 2004.

\bibitem{Wood1991}
J. C. Wood,
\newblock Lewy's theorem fails in higher dimensions,
\newblock \emph{Mathematica Scandinavica} \textbf{69} (1991), no. 2, 166--170.
\newblock DOI: \texttt{10.7146/math.scand.a-12375}.

\end{thebibliography}

\end{document}







